% Start of user code protected header
\documentclass{gemoc} %% no option needed, default is : 10pt, twoside, babel[ english] , graphicx
\usepackage{color}
\usepackage[colorlinks=true]{hyperref}

\task{x.x.x}
\title{Gemoc XXX }
\docnumber{Dx.x.x}
\version{1.0}

\companycopyright{Consortium GEMOC}  %% Appear in the foot page

\begin{document}
\maketitle

\begin{revisions}
	\begin{revtable}
		\dates{}{}{}{}{}
		\writers{}{}{}{}{}
		\approvers{}{}{}{}{}
	\end{revtable}
	\begin{revisionlabels}
		\revlabel{}
	\end{revisionlabels}
\end{revisions}
\begin{tableofauthors}
	\leadauthor{Didier Vojtisek}{INRIA}
	\contributor{$<$Name$>$}{$<$Organisation$>$}
\end{tableofauthors}

\tableofcontents
\newpage

\chapter{Gemoc language workbench workflow}

%End of user code
%Start of user code summary
The figure \ref{fig:Gemoc_language_workbench_workflow} presents the global view of the workflow of the different activities of the Language Workbench.


In further sections each activity will be detailled by also presenting the major concrete artefacts resulting from the Commands.
%End of user code
\begin{figure}[h!]
		\center
		\includegraphics*[trim=0.0cm 0.0cm 0cm 0.0cm, clip=true]{fig/Gemoc_language_workbench_workflow}
		\caption{Gemoc Language Workbench workflow activities and supported tools}
		\label{fig:Gemoc_language_workbench_workflow}
\end{figure}

% Start of user code protected summary_legend
Legend:
\begin{itemize}
	\item \includegraphics[width=1cm]{fig/step} Activity: an activity is a step of the workflow. Each activity is supported by at least one concrete Command
	\item \includegraphics[width=0.7cm]{fig/command} Command: a command is a concrete action of the studio, usually implemented as a wizard.
	\item \includegraphics[width=0.5cm]{fig/artifact_add} Artifact creation: Artifact created as the result of a command.
	\item \includegraphics[width=0.5cm]{fig/artifact_update} Artifact update:  Artifact updated as the result of a command.
\end{itemize}
%End of user code
\section{Create language workbench definition}
%%%%%%%%%%%%%%%%%%%%%%%%%%%%%%%%%%%%%%%%%%%%%%%%%%%%%%%
\label{sec:Create_language_workbench_definition}
% Start of user code protected Create language workbench definition
%End of user code
\begin{figure}[h!]
		\center
		\includegraphics*[trim=0.0cm 0.0cm 0cm 0.0cm, clip=true]{fig/Create_language_workbench_definition}
		\caption{Create language workbench definition activity}
		\label{fig:Create_language_workbench_definition}
\end{figure}

TODO

The figure \ref{fig:Create_language_workbench_definition} presents the activity and its supporting Commands.

\subsection{Create language workbench Command}
TODO
\subsubsection{Created artefacts}
Artifacts created by the Create language workbench Command:
\paragraph{Language workbench project} 
TODO\paragraph{xdsml file} 
TODO
\subsubsection{Updated artefacts}
Artifacts updated by the Create language workbench Command:

	None

\section{Define language}
%%%%%%%%%%%%%%%%%%%%%%%%%%%%%%%%%%%%%%%%%%%%%%%%%%%%%%%
\label{sec:Define_language}
% Start of user code protected Define language
%End of user code
\begin{figure}[h!]
		\center
		\includegraphics*[trim=0.0cm 0.0cm 0cm 0.0cm, clip=true]{fig/Define_language}
		\caption{Define language activity}
		\label{fig:Define_language}
\end{figure}

TODO

The figure \ref{fig:Define_language} presents the activity and its supporting Commands.

\subsection{Create EMF Project Command}
TODO
\subsubsection{Created artefacts}
Artifacts created by the Create EMF Project Command:
\paragraph{EMF project} 
TODO\paragraph{ecore} 
TODO\paragraph{genmodel} 
TODO
\subsubsection{Updated artefacts}
Artifacts updated by the Create EMF Project Command:
\paragraph{xdsml} 
TODO

\section{Define editors}
%%%%%%%%%%%%%%%%%%%%%%%%%%%%%%%%%%%%%%%%%%%%%%%%%%%%%%%
\label{sec:Define_editors}
% Start of user code protected Define editors
%End of user code
\begin{figure}[h!]
		\center
		\includegraphics*[trim=0.0cm 0.0cm 0cm 0.0cm, clip=true]{fig/Define_editors}
		\caption{Define editors activity}
		\label{fig:Define_editors}
\end{figure}

TODO

The figure \ref{fig:Define_editors} presents the activity and its supporting Commands.

\subsection{Create tree editor Command}
TODO
\subsubsection{Created artefacts}
Artifacts created by the Create tree editor Command:
\paragraph{Edit project} 
TODO\paragraph{Editor project} 
TODO
\subsubsection{Updated artefacts}
Artifacts updated by the Create tree editor Command:
\paragraph{xdsml} 
TODO

\subsection{Create Xtext editor Command}
TODO
\subsubsection{Created artefacts}
Artifacts created by the Create Xtext editor Command:
\paragraph{Xtext project} 
TODO\paragraph{Xtext.ui project} 
TODO
\subsubsection{Updated artefacts}
Artifacts updated by the Create Xtext editor Command:
\paragraph{xdsml} 
TODO

\subsection{Create Sirius editor Command}
TODO
\subsubsection{Created artefacts}
Artifacts created by the Create Sirius editor Command:
\paragraph{Sirius specification project} 
TODO
\subsubsection{Updated artefacts}
Artifacts updated by the Create Sirius editor Command:
\paragraph{xdsml} 
TODO

\section{Define DSA}
%%%%%%%%%%%%%%%%%%%%%%%%%%%%%%%%%%%%%%%%%%%%%%%%%%%%%%%
\label{sec:Define_DSA}
% Start of user code protected Define DSA
%End of user code
\begin{figure}[h!]
		\center
		\includegraphics*[trim=0.0cm 0.0cm 0cm 0.0cm, clip=true]{fig/Define_DSA}
		\caption{Define DSA activity}
		\label{fig:Define_DSA}
\end{figure}

TODO

The figure \ref{fig:Define_DSA} presents the activity and its supporting Commands.

\subsection{K2 Command}
TODO
\subsubsection{Created artefacts}
Artifacts created by the K2 Command:
\paragraph{K2 project} 
TODO
\subsubsection{Updated artefacts}
Artifacts updated by the K2 Command:
\paragraph{xdsml} 
TODO

\subsection{K3 Command}
TODO
\subsubsection{Created artefacts}
Artifacts created by the K3 Command:
\paragraph{K3 project} 
TODO
\subsubsection{Updated artefacts}
Artifacts updated by the K3 Command:
\paragraph{xdsml} 
TODO

\section{Define MOC}
%%%%%%%%%%%%%%%%%%%%%%%%%%%%%%%%%%%%%%%%%%%%%%%%%%%%%%%
\label{sec:Define_MOC}
% Start of user code protected Define MOC
%End of user code
\begin{figure}[h!]
		\center
		\includegraphics*[trim=0.0cm 0.0cm 0cm 0.0cm, clip=true]{fig/Define_MOC}
		\caption{Define MOC activity}
		\label{fig:Define_MOC}
\end{figure}

TODO

The figure \ref{fig:Define_MOC} presents the activity and its supporting Commands.

\subsection{Moc project Command}
TODO
\subsubsection{Created artefacts}
Artifacts created by the Moc project Command:

	None
\subsubsection{Updated artefacts}
Artifacts updated by the Moc project Command:

	None

\section{DSE project}
%%%%%%%%%%%%%%%%%%%%%%%%%%%%%%%%%%%%%%%%%%%%%%%%%%%%%%%
\label{sec:DSE_project}
% Start of user code protected DSE project
%End of user code
\begin{figure}[h!]
		\center
		\includegraphics*[trim=0.0cm 0.0cm 0cm 0.0cm, clip=true]{fig/DSE_project}
		\caption{DSE project activity}
		\label{fig:DSE_project}
\end{figure}

TODO

The figure \ref{fig:DSE_project} presents the activity and its supporting Commands.

\subsection{ECL Command}
TODO
\subsubsection{Created artefacts}
Artifacts created by the ECL Command:
\paragraph{ecl file} 
TODO
\subsubsection{Updated artefacts}
Artifacts updated by the ECL Command:
\paragraph{xdsml} 
TODO

\subsection{Modelhex Command}
TODO
\subsubsection{Created artefacts}
Artifacts created by the Modelhex Command:

	None
\subsubsection{Updated artefacts}
Artifacts updated by the Modelhex Command:

	None

\section{Animator}
%%%%%%%%%%%%%%%%%%%%%%%%%%%%%%%%%%%%%%%%%%%%%%%%%%%%%%%
\label{sec:Animator}
% Start of user code protected Animator
%End of user code
\begin{figure}[h!]
		\center
		\includegraphics*[trim=0.0cm 0.0cm 0cm 0.0cm, clip=true]{fig/Animator}
		\caption{Animator activity}
		\label{fig:Animator}
\end{figure}

TODO

The figure \ref{fig:Animator} presents the activity and its supporting Commands.


%Start of user code protected footer
\end{document}
%End of user code
